% SOURCE_AUTHORITY: sga5_fr_workpass.tex, lines 5758--5835 (LF-normalized unit).
% SOURCE_SHA256: 791F4EFFC5E02832D5D77ED03518C8156D6F07E4C8238B03545DB93D883FBB28
\subsection*{6.6}
Sejam $A$ uma $\Lambda$-álgebra plana, e, para $i=1,2$, $X_i$ um $S$-esquema, $X=X_1\times_SX_2$, $L_i\in\Ob D_{\mathrm{ctf}}({}_AX_i)$. Define-se um emparelhamento
\begin{equation}
\uRHom_A(p_1^*L_1,p_2^!L_2)
\lten_{\Lambda}
\uRHom_A(p_2^*L_2,p_1^!L_1)
\longrightarrow
KA_{X,L_{\natural}}
\tag{6.6.1}
\end{equation}
da seguinte forma. Identifica-se o primeiro membro com
\[
\uRHom_A(p_1^*L_1,p_1^*KA_{X_1})
\lten_A p_2^*L_2
\lten_{\Lambda}
\uRHom_A(p_2^*L_2,p_2^*KA_{X_2})
\lten_A p_1^*L_1
\]
por meio dos isomorfismos 6.4.1, e depois com
\[
A\lten_{A^e}\uRHom_A(\ )\lten_{\Lambda}p_2^*L_2
 \lten_{\Lambda}\uRHom_A(\ )\lten_{\Lambda}p_1^*L_1\lten_{A^e}A
\]
por meio de 5.3.1; este objeto é enviado a
\[
A\lten_{A^e}(KA_{X_1}\lten_{\Lambda}KA_{X_2})\lten_{A^e}A=R
\]
pelo produto tensorial das flechas de avaliação 5.3.6; por último, identifica-se $R$ com
\[
A\lten_{A^e}(KA_{X_1}\lten_A KA_{X_2})
\]
por meio de 5.8.2, e depois com $KA_{X,L_{\natural}}$ por meio de 6.5.2 e (III 1.7.6).

Se $c:C\to X$, $d:D\to X$ são correspondências, e
\[
e=c\times_Xd:E\to X,
\]
de 6.6.1 deduz-se, como em (III 4.2), um emparelhamento
\begin{equation}
c_*\uRHom_A(c_1^*L_1,c_2^!L_2)
\lten_{\Lambda}
d_*\uRHom_A(d_2^*L_2,d_1^!L_1)
\longrightarrow
e_*KA_{E,L_{\natural}},
\tag{6.6.2}
\end{equation}
de onde se obtém um produto cup
\begin{equation}
\langle\ ,\ \rangle_A:
\uHom_A(c_1^*L_1,c_2^!L_2)\otimes
\uHom_A(d_2^*L_2,d_1^!L_1)
\longrightarrow
H^0(E,KA_{E,L_{\natural}}).
\tag{6.6.3}
\end{equation}
Não é difícil estender ao caso não comutativo a noção de composição de correspondências estudada em (III 5.2), e verifica-se que, para
\[
u\in\uHom_A(c_1^*L_1,c_2^!L_2),\qquad
v\in\uHom_A(d_2^*L_2,d_1^!L_1),
\]
tem-se a fórmula análoga a (III 5.2.9)
\begin{equation}
\langle u,v\rangle_A=\Tr_A(vu)=\Tr_A(uv).
\tag{6.6.4}
\end{equation}
A fórmula de Lefschetz (III 4.4) generaliza-se sem dificuldade neste contexto: dados $S$-morfismos próprios $f_i:X_i\to Y_i$ $(i=1,2)$, e um cubo comutativo (III 4.4.0), tem-se, para $L_i\in\Ob D_{\mathrm{ctf}}({}_AX_i)$, um quadrado comutativo análogo a (III 4.4.1), com $K_C$ (resp. $K_D$) substituído por $KA_{C,L_{\natural}}$ (resp. $KA_{D,L_{\natural}}$), e uma fórmula análoga a (III 4.5.1), em $H^0(D,KA_{D,L_{\natural}})$:
\begin{equation}
\langle f_*u,f_*v\rangle_A=g_*\langle u,v\rangle_A.
\tag{6.6.5}
\end{equation}
A flecha
\[
f_*c_*KA_{C,L_{\natural}}\longrightarrow d_*KA_{D,L_{\natural}}
\]
deduz-se da flecha (4) de (III 4.4.1) por tensorização com
\[
A\lten_{A^e}A .
\]
